\chapter{Metodologia}\label{cap:tres}

Este capítulo apresenta a metodologia proposta para estimação de fasores harmônicos em condições de frequência variável. A proposta é organizada como uma cadeia de processamento em tempo discreto, composta por estimação da frequência fundamental, reamostragem dinâmica por interpolação B-Spline e extração harmônica por banco de filtros Flat-Top em estrutura polifásica. O objetivo central é produzir uma sequência de amostras sincronizada com a frequência instantânea do sinal, de modo que o banco de filtros opere em uma base temporal mais adequada à estimação dos fasores harmônicos.

A metodologia descrita neste capítulo corresponde ao algoritmo que orienta a implementação em desenvolvimento. Assim, as equações e etapas apresentadas têm caráter metodológico: definem o processamento necessário, os sinais de entrada e saída de cada bloco e as métricas que serão utilizadas na avaliação posterior. Os resultados finais de desempenho e a validação completa em FPGA não são tratados neste capítulo.

% Figura futura recomendada:
% Entrada -> pré-filtro ZC -> Zero-Crossing/interpolação -> média móvel -> atraso/alinhamento -> pré-filtro B-Spline -> Farrow -> banco Flat-Top polifásico -> fasores harmônicos -> métricas.
% Esta figura deve ser refeita em TikZ ou ferramenta vetorial limpa, sem anotações de desenvolvimento.

\section{Visão geral do método proposto}

O método parte de um sinal discreto $x[n]$, obtido a partir da amostragem de um sinal elétrico com frequência fundamental possivelmente diferente da frequência nominal. Em aplicações de estimação fasorial, esse desvio de frequência é crítico, pois a maior parte dos estimadores baseados em transformadas ou bancos de filtros assume uma relação bem definida entre frequência fundamental, janela de observação e pontos espectrais avaliados. Quando essa relação é perdida, ocorre desalinhamento temporal e espectral, o que afeta principalmente a fase dos fasores estimados e se torna mais severo em componentes harmônicas de ordem elevada.

Para contornar esse problema, a metodologia proposta emprega uma etapa de rastreamento de frequência por cruzamento por zero, seguida por uma etapa de reamostragem dinâmica. A frequência estimada controla um interpolador B-Spline implementado por estrutura de Farrow, que produz uma sequência reamostrada $x_i[k]$ com número de amostras por ciclo associado à frequência nominal. Em seguida, essa sequência alimenta um banco de filtros Flat-Top polifásico, responsável pela extração simultânea dos fasores da componente fundamental e das componentes harmônicas de interesse.

De forma resumida, a cadeia de processamento é composta pelas seguintes etapas:

\begin{itemize}
    \item geração ou aquisição do sinal de entrada com componentes harmônicas;
    \item filtragem preliminar para robustecer a detecção de cruzamentos por zero;
    \item estimação da frequência fundamental por Zero-Crossing adaptado;
    \item suavização da frequência estimada por média móvel;
    \item alinhamento temporal entre sinal e frequência estimada;
    \item reamostragem por interpolação B-Spline cúbica;
    \item extração harmônica por banco de filtros Flat-Top com decomposição polifásica;
    \item correção de fase e cálculo das métricas de avaliação.
\end{itemize}

A escolha desses blocos está relacionada ao compromisso entre precisão, robustez e custo computacional. O estimador por cruzamento por zero é simples e adequado a operação em tempo real, desde que precedido de filtragem apropriada. A interpolação B-Spline permite deslocar o problema do reprojeto contínuo de filtros para uma etapa de reamostragem controlada pela frequência estimada \cite{aleixo2022,aleixoTese}. Por fim, a decomposição polifásica reduz o esforço necessário para implementar o banco de filtros, preservando a estrutura de extração harmônica baseada em filtros Flat-Top \cite{duda2019,mk06}.

\section{Arquitetura funcional da metodologia}

A metodologia pode ser organizada em três níveis funcionais: definição dos sinais de teste, sincronização temporal por estimação de frequência e reamostragem, e extração dos fasores harmônicos. Essa organização não representa uma arquitetura final de hardware, mas uma descrição algorítmica dos blocos necessários para que o estimador opere com sinais cuja frequência fundamental pode se afastar da nominal.

No primeiro nível, são definidos os sinais de entrada e as condições de avaliação. O sinal $x[n]$ pode conter componente fundamental, componentes harmônicas, ruído, frequência \textit{off-nominal}, rampa de frequência, modulação de amplitude e modulação de fase. Esses cenários são tratados como condições de validação inspiradas na IEC/IEEE 60255-118-1, adaptadas ao estudo de fasores harmônicos. Também são definidos parâmetros como frequência nominal, número de pontos por ciclo, relação sinal-ruído, amplitudes harmônicas, ordens harmônicas avaliadas e resolução numérica.

No segundo nível, encontram-se as etapas responsáveis por estimar a frequência fundamental e ajustar a base temporal do sinal. Inicialmente, o sinal de entrada passa por um pré-filtro associado ao Zero-Crossing, com o objetivo de atenuar componentes harmônicas e ruído antes da detecção dos cruzamentos por zero. Em seguida, o instante de cruzamento é refinado por interpolação linear, reduzindo o erro causado pela amostragem discreta. A estimativa é atualizada a cada meio ciclo, o que melhora o acompanhamento de variações de frequência em ensaios dinâmicos, como rampas.

A frequência obtida pelo Zero-Crossing é suavizada por média móvel e alinhada temporalmente com o trecho do sinal ao qual ela se refere. Esse alinhamento é necessário porque o pré-filtro do Zero-Crossing, a média móvel e o pré-filtro B-Spline introduzem atrasos. Se esses atrasos não forem compensados, o interpolador pode usar uma frequência estimada para um trecho diferente daquele que está sendo reamostrado, produzindo erro de fase. Após a compensação, o sinal é processado pelo pré-filtro B-Spline e pela estrutura de Farrow, gerando a sequência reamostrada $x_i[k]$.

No terceiro nível, a sequência reamostrada alimenta o banco de filtros harmônicos. O filtro Flat-Top é usado como protótipo por apresentar baixa sensibilidade de magnitude a pequenos desvios em torno da frequência central. Para reduzir o custo computacional, seus coeficientes são reorganizados por decomposição polifásica, e as saídas dos ramos são combinadas por uma IFFT para produzir simultaneamente os canais associados às componentes de interesse. A partir desses canais, obtêm-se os fasores harmônicos estimados.

As métricas de avaliação são calculadas apenas após a extração fasorial. Para cada componente, os fasores estimados serão comparados com fasores de referência por meio de TVE, erro de magnitude, erro de fase, erro de frequência e erro de ROCOF. Dessa forma, a metodologia separa claramente as etapas de geração do sinal, sincronização e reamostragem, extração harmônica e avaliação, o que facilita a análise posterior dos erros introduzidos por cada bloco.

\section{Modelo dos sinais de teste}

Embora os resultados finais sejam discutidos apenas posteriormente, a metodologia de validação precisa ser definida junto com o algoritmo. Os sinais utilizados para análise são inspirados nos ensaios da norma IEC/IEEE 60255-118-1, que estabelece critérios para PMUs na estimação da componente fundamental, frequência e ROCOF \cite{iec60255}. Como ainda não há uma norma equivalente consolidada para fasores harmônicos, os ensaios normativos são utilizados como referência metodológica para avaliar o comportamento do estimador quando o sinal contém componentes harmônicas.

De forma geral, considera-se um sinal composto pela componente fundamental e por componentes harmônicas, escrito como

\begin{equation}
x(t) = A_1(t)\cos\left(\theta_1(t)\right) + \sum_{h \in \mathcal{H}} A_h(t)\cos\left(h\theta_1(t) + \varphi_h\right) + \eta(t),
\label{eq:sinal_harmonico_metodologia}
\end{equation}

\noindent em que $A_1(t)$ é a amplitude da componente fundamental, $A_h(t)$ é a amplitude da componente harmônica de ordem $h$, $\theta_1(t)$ é a fase instantânea da fundamental, $\varphi_h$ é a fase inicial da componente harmônica e $\eta(t)$ representa ruído aditivo. O conjunto $\mathcal{H}$ contém as ordens harmônicas avaliadas. Nos ensaios considerados neste trabalho, o interesse principal recai sobre componentes ímpares até a 49ª ordem, coerentemente com a aplicação de monitoramento harmônico em sistemas elétricos.

A frequência instantânea da fundamental é obtida a partir da derivada da fase,

\begin{equation}
f(t) = \frac{1}{2\pi}\frac{d\theta_1(t)}{dt}.
\label{eq:freq_instantanea}
\end{equation}

Assim, diferentes cenários de teste podem ser descritos por diferentes escolhas de $A_1(t)$, $A_h(t)$ e $f(t)$. Em condição \textit{off-nominal}, a frequência é constante, porém diferente da nominal. Em rampa de frequência, $f(t)$ varia linearmente no tempo. Nos ensaios de modulação de amplitude e modulação de fase, a amplitude ou a fase são perturbadas por sinais senoidais de baixa frequência. Essas condições permitem avaliar se a cadeia de estimação mantém coerência temporal quando o sinal deixa de ser estritamente estacionário.

No domínio discreto, o sinal amostrado é representado por

\begin{equation}
x[n] = x(nT_s),
\label{eq:sinal_discreto}
\end{equation}

\noindent em que $T_s = 1/F_s$ é o período de amostragem. Na implementação de referência em software, utiliza-se $F_s = f_0N_{ppc}$, sendo $f_0$ a frequência nominal e $N_{ppc}$ o número nominal de pontos por ciclo. Essa escolha facilita a conexão entre o interpolador e o banco de filtros, pois o fator de decimação do banco polifásico pode ser associado ao número nominal de amostras por ciclo.

\section{Estimação de frequência por Zero-Crossing}

A primeira etapa dinâmica do método é a estimação da frequência fundamental. Para isso, emprega-se um detector de cruzamento por zero, escolhido por sua simplicidade e por sua baixa demanda computacional \cite{ribeiro2013}. Antes da detecção, o sinal é processado por um filtro voltado à atenuação de harmônicos, inter-harmônicos e ruído. Essa etapa é importante porque componentes de alta frequência podem introduzir cruzamentos espúrios, fazendo com que o período estimado deixe de representar a componente fundamental.

Considere duas amostras consecutivas do sinal filtrado, $x_f[n-1]$ e $x_f[n]$, entre as quais ocorre mudança de sinal. O instante de cruzamento pode ser refinado por interpolação linear. Se $t_{n-1}$ é o instante da amostra anterior, uma estimativa do instante de cruzamento é dada por

\begin{equation}
\hat{t}_{zc} = t_{n-1} + T_s \frac{x_f[n-1]}{x_f[n-1] - x_f[n]}.
\label{eq:interpolacao_zc}
\end{equation}

O termo fracionário da Equação~\eqref{eq:interpolacao_zc} reduz o erro de quantização temporal que ocorreria caso o cruzamento fosse associado diretamente ao índice da amostra. Esse refinamento é especialmente importante quando a frequência estimada será usada para controlar a reamostragem, pois pequenos erros de tempo podem se acumular na fase das componentes harmônicas.

No algoritmo adotado, a frequência é estimada a partir do intervalo entre cruzamentos separados por meio ciclo. Essa escolha permite atualizar a estimativa com maior frequência do que uma estratégia baseada no ciclo completo, o que é útil em cenários de rampa, nos quais a frequência varia durante o próprio ciclo. Se $\Delta \hat{t}_{zc}$ representa o intervalo estimado entre dois cruzamentos consecutivos de polaridade oposta, a frequência estimada é

\begin{equation}
\hat{f}_{zc}[n] = \frac{1}{2\Delta \hat{t}_{zc}}.
\label{eq:freq_zc}
\end{equation}

A Equação~\eqref{eq:freq_zc} mostra que o estimador depende diretamente da precisão dos instantes de cruzamento. Por isso, a etapa de filtragem anterior ao Zero-Crossing e a interpolação linear do cruzamento são tratadas como partes essenciais do método, e não apenas como detalhes de implementação.

\section{Suavização e alinhamento temporal}

A sequência $\hat{f}_{zc}[n]$ obtida pelo detector de cruzamento por zero pode apresentar oscilações decorrentes de ruído, distorções residuais e erros numéricos. Para reduzir essas variações, aplica-se uma média móvel de comprimento $N_w$, definida por

\begin{equation}
\hat{f}[n] = \frac{1}{N_w}\sum_{i=0}^{N_w-1}\hat{f}_{zc}[n-i].
\label{eq:media_movel_freq}
\end{equation}

Essa filtragem suaviza a frequência utilizada pelo interpolador, reduzindo variações bruscas no parâmetro de reamostragem. Como consequência, ela também introduz atraso de grupo. Além disso, o pré-filtro usado antes do Zero-Crossing e o pré-filtro B-Spline também introduzem atrasos próprios. Para que a amostra do sinal e a frequência usada na reamostragem correspondam ao mesmo trecho temporal do fenômeno, é necessário alinhar as sequências por meio de atrasos explícitos.

O sinal alinhado pode ser representado por

\begin{equation}
x_d[n] = x[n-D_x],
\label{eq:atraso_sinal}
\end{equation}

\noindent enquanto a frequência suavizada e alinhada é representada por

\begin{equation}
\hat{f}_d[n] = \hat{f}[n-D_f].
\label{eq:atraso_freq}
\end{equation}

Os atrasos $D_x$ e $D_f$ não representam uma alteração física do sinal, mas uma compensação necessária para manter coerência entre os blocos. Em uma implementação em tempo real, esse alinhamento é feito por buffers. Na implementação de referência, o atraso total é escolhido a partir dos atrasos do filtro anterior ao Zero-Crossing, da média móvel e do pré-filtro associado ao interpolador.

\section{Reamostragem por interpolação B-Spline}

Após a estimação da frequência, o sinal é reamostrado para reduzir o desalinhamento entre a frequência real do sinal e a base de estimação do banco de filtros. A ideia é produzir uma sequência $x_i[k]$ que preserve o número nominal de pontos por ciclo associado a $f_0$, mesmo quando a frequência real se desvia da nominal. Para isso, define-se a razão instantânea de reamostragem

\begin{equation}
\lambda[n] = \frac{f_0}{\hat{f}_d[n]}.
\label{eq:lambda_reamostragem}
\end{equation}

Quando $\hat{f}_d[n] > f_0$, a razão $\lambda[n]$ torna-se menor que um; quando $\hat{f}_d[n] < f_0$, torna-se maior que um. Dessa forma, a taxa efetiva de geração das amostras interpoladas acompanha a frequência estimada, evitando que o banco de filtros opere sobre uma sequência temporal incompatível com sua base nominal.

O interpolador utilizado é baseado em B-Splines cúbicas com pré-filtragem, conforme discutido em \citeonline{aleixo2022}. A interpolação B-Spline possui boas propriedades de suavidade, mas a forma cúbica básica não passa necessariamente pelos valores originais das amostras. Por isso, emprega-se um pré-filtro que transforma a sequência de amostras em coeficientes adequados para a reconstrução interpolante.

No caso cúbico, o pré-filtro pode ser representado por uma função de transferência racional,

\begin{equation}
H_{pre}(z) = \frac{B(z)}{A(z)},
\label{eq:prefiltro_bspline}
\end{equation}

\noindent de modo que a sequência pré-filtrada $c[n]$ seja obtida por

\begin{equation}
c[n] = H_{pre}(z)\{x_d[n]\}.
\label{eq:coef_bspline}
\end{equation}

Na implementação adotada, os coeficientes desse pré-filtro são calculados a partir do polo associado à B-Spline cúbica, frequentemente escrito como

\begin{equation}
s = \sqrt{3} - 2.
\label{eq:polo_bspline}
\end{equation}

O pré-filtro compensa a suavização inerente à base B-Spline e melhora a fidelidade da amostra interpolada. Essa etapa é particularmente relevante para harmônicos de ordem elevada, pois a atenuação ou o atraso de fase introduzidos na reconstrução temporal podem produzir erros crescentes com a ordem harmônica.

\section{Estrutura de Farrow}

A geração de amostras em instantes fracionários é realizada por uma estrutura de Farrow. Essa estrutura implementa um filtro de atraso fracionário variável como um polinômio em função do parâmetro fracionário $\alpha$. Para a interpolação cúbica utilizada neste trabalho, a amostra interpolada pode ser escrita como

\begin{equation}
x_i[k] = H_0[n]\alpha^3 + H_1[n]\alpha^2 + H_2[n]\alpha + H_3[n],
\label{eq:farrow_polinomio}
\end{equation}

\noindent em que $\alpha \in [0,1)$ indica a posição fracionária da amostra desejada em relação às amostras disponíveis, e $H_0[n]$, $H_1[n]$, $H_2[n]$ e $H_3[n]$ são combinações lineares de quatro amostras consecutivas do sinal pré-filtrado. Considerando o vetor local

\begin{equation}
\mathbf{c}_n =
\begin{bmatrix}
c[n-3] & c[n-2] & c[n-1] & c[n]
\end{bmatrix}^{T},
\label{eq:vetor_farrow}
\end{equation}

\noindent os subfiltros cúbicos podem ser expressos como

\begin{align}
H_0[n] &= -\frac{1}{6}c[n-3] + \frac{1}{2}c[n-2] - \frac{1}{2}c[n-1] + \frac{1}{6}c[n],
\label{eq:farrow_h0}\\
H_1[n] &= \frac{1}{2}c[n-3] - c[n-2] + \frac{1}{2}c[n-1],
\label{eq:farrow_h1}\\
H_2[n] &= -\frac{1}{2}c[n-3] + \frac{1}{2}c[n-1],
\label{eq:farrow_h2}\\
H_3[n] &= \frac{1}{6}c[n-3] + \frac{2}{3}c[n-2] + \frac{1}{6}c[n-1].
\label{eq:farrow_h3}
\end{align}

O parâmetro $\alpha$ é atualizado de acordo com a razão $\lambda[n]$. Enquanto $\alpha < 1$, novas amostras interpoladas são geradas. Quando $\alpha$ ultrapassa a unidade, subtrai-se um período de amostra para reposicionar o acumulador fracionário:

\begin{equation}
\alpha \leftarrow \alpha + \lambda[n],
\qquad
\alpha \leftarrow \alpha - 1 \quad \text{quando } \alpha \geq 1.
\label{eq:atualizacao_alpha}
\end{equation}

Essa forma de implementação evita o recálculo dos coeficientes do filtro para cada novo atraso fracionário. Os coeficientes dos subfiltros permanecem fixos, e apenas as potências de $\alpha$ mudam ao longo do tempo. Essa característica torna a estrutura adequada para aplicações embarcadas e para implementação em arquiteturas de processamento com recursos limitados.

\section{Banco de filtros Flat-Top}

Após a reamostragem, a sequência $x_i[k]$ alimenta o banco de filtros responsável pela estimação dos fasores harmônicos. O filtro protótipo adotado é do tipo Flat-Top, cuja principal característica é apresentar resposta de magnitude bastante plana em torno da frequência central. Essa propriedade é desejável para estimação de amplitude, pois reduz a sensibilidade a pequenos desvios de frequência dentro da banda de passagem \cite{duda2016,duda2019}.

O filtro protótipo pode ser definido por uma janela cossenoidal Flat-Top normalizada. Uma forma geral para seus coeficientes é

\begin{equation}
w[n] = \frac{1}{C}\sum_{m=0}^{P} a_m \cos\left(\frac{2\pi mn}{N}\right),
\label{eq:janela_flattop}
\end{equation}

\noindent em que $a_m$ são os coeficientes da janela, $N$ é o comprimento do filtro, $P$ é a ordem da combinação cossenoidal e $C$ é um fator de normalização escolhido para que o ganho em corrente contínua seja unitário. A partir desse filtro base, o filtro associado à componente harmônica de ordem $h$ é obtido por modulação complexa:

\begin{equation}
h_h[n] = w[n]e^{j2\pi h f_0 n/F_s}.
\label{eq:filtro_harmonico}
\end{equation}

A saída do filtro associado à ordem $h$ é então

\begin{equation}
X_h[k] = \sum_{n=0}^{N-1} h_h[n]x_i[k-n].
\label{eq:saida_filtro_harmonico}
\end{equation}

A implementação direta da Equação~\eqref{eq:saida_filtro_harmonico} para muitas ordens harmônicas pode se tornar custosa. Se o banco possuir $M$ canais e o filtro protótipo possuir $N$ coeficientes, a forma direta exige uma quantidade de operações proporcional a $NM$ para produzir uma saída completa do banco. Por esse motivo, utiliza-se a estrutura polifásica.

\section{Decomposição polifásica}

A decomposição polifásica permite particionar o filtro protótipo em $M$ componentes, cada um operando em uma taxa reduzida. Para um filtro $H(z)$, sua decomposição polifásica do Tipo I pode ser escrita como \cite{mk06}

\begin{equation}
H(z) = \sum_{\ell=0}^{M-1} z^{-\ell}E_{\ell}(z^M),
\label{eq:decomposicao_polifasica}
\end{equation}

\noindent em que

\begin{equation}
E_{\ell}(z) = \sum_{r=-\infty}^{\infty} h[rM+\ell]z^{-r},
\qquad 0 \leq \ell \leq M-1.
\label{eq:componentes_polifasicos}
\end{equation}

Essa decomposição reorganiza os coeficientes do filtro em ramos distintos. Quando aplicada ao banco uniforme, as saídas dos componentes polifásicos podem ser combinadas por uma transformada discreta inversa, produzindo simultaneamente as saídas dos canais do banco. Em termos práticos, a estrutura polifásica desloca parte do processamento para uma taxa decimada e substitui a implementação direta de vários filtros modulados por uma combinação entre filtragem polifásica e IFFT.

Para um banco com $M$ canais e filtro protótipo de comprimento $N$, a complexidade aproximada passa de $NM$ multiplicações, na forma direta, para

\begin{equation}
N + M\log_2(M),
\label{eq:complexidade_polifasica}
\end{equation}

\noindent quando a etapa de modulação é realizada por uma FFT ou IFFT. Essa redução é uma das principais motivações para o uso da decomposição polifásica no presente trabalho, pois a estimação simultânea de vários fasores harmônicos exige um banco com muitos canais.

\section{Correção de fase}

Mesmo após a reamostragem, a saída do banco de filtros pode necessitar de correção de fase associada ao desvio acumulado entre a frequência estimada e a frequência nominal. Seja

\begin{equation}
\Delta f[k] = \hat{f}[k] - f_0.
\label{eq:desvio_freq}
\end{equation}

Uma aproximação discreta da fase acumulada devido ao desvio de frequência pode ser obtida por integração numérica. Utilizando a regra trapezoidal, define-se

\begin{equation}
\Phi[k] = \Phi[k-1] + \pi\left(\Delta f[k] + \Delta f[k-1]\right)T_r,
\label{eq:correcao_fase_fundamental}
\end{equation}

\noindent em que $T_r$ é o período efetivo entre saídas do banco após a decimação. Para a componente harmônica de ordem $h$, a correção é proporcional à ordem harmônica:

\begin{equation}
\Phi_h[k] = h\Phi[k].
\label{eq:correcao_fase_harmonico}
\end{equation}

Assim, o fasor corrigido é calculado a partir da magnitude estimada e da fase ajustada,

\begin{equation}
\hat{X}_h[k] = |\hat{X}_h[k]|e^{j(\angle \hat{X}_h[k] + \Phi_h[k])}.
\label{eq:fasor_corrigido}
\end{equation}

Essa etapa é importante porque pequenos desvios de fase na fundamental são amplificados nas ordens harmônicas. Portanto, erros de alinhamento temporal, quantização da frequência ou atrasos não compensados tendem a aparecer de forma mais evidente nos fasores harmônicos de maior ordem.

\section{Métricas para avaliação}

A avaliação do método será baseada em métricas usuais de estimação fasorial. A principal delas é o Erro Vetorial Total, ou \textit{Total Vector Error} (TVE), definido como

\begin{equation}
TVE_h[k] =
100
\sqrt{
\frac{
\left(\operatorname{Re}\{\hat{X}_h[k]\}-\operatorname{Re}\{X_h[k]\}\right)^2+
\left(\operatorname{Im}\{\hat{X}_h[k]\}-\operatorname{Im}\{X_h[k]\}\right)^2
}{
\left(\operatorname{Re}\{X_h[k]\}\right)^2+
\left(\operatorname{Im}\{X_h[k]\}\right)^2
}
},
\label{eq:tve_metodologia}
\end{equation}

\noindent em que $\hat{X}_h[k]$ é o fasor estimado e $X_h[k]$ é o fasor de referência da componente de ordem $h$. O TVE combina erro de magnitude e erro de fase em uma única medida vetorial, o que o torna adequado para comparar o fasor estimado com a referência.

Também são considerados o erro de magnitude e o erro de fase, definidos por

\begin{equation}
\varepsilon_{A,h}[k] = |\hat{X}_h[k]| - |X_h[k]|,
\label{eq:erro_magnitude}
\end{equation}

\begin{equation}
\varepsilon_{\varphi,h}[k] = \angle \hat{X}_h[k] - \angle X_h[k].
\label{eq:erro_fase}
\end{equation}

Essas duas métricas permitem interpretar a origem do TVE. Em sinais com variação de frequência, por exemplo, é comum que erros de alinhamento temporal se manifestem principalmente como erro de fase, e que esse efeito cresça com a ordem harmônica. Por fim, quando a frequência estimada também é avaliada, pode-se utilizar o erro de frequência

\begin{equation}
FE[k] = \hat{f}[k] - f[k],
\label{eq:erro_frequencia}
\end{equation}

\noindent e o erro de taxa de variação da frequência, associado ao ROCOF,

\begin{equation}
RFE[k] = \widehat{ROCOF}[k] - ROCOF[k].
\label{eq:erro_rocof}
\end{equation}

No contexto deste trabalho, essas métricas serão utilizadas para organizar a validação posterior do método em diferentes cenários de teste. Como os resultados finais ainda estão em desenvolvimento, este capítulo limita-se a estabelecer as definições e o encadeamento metodológico necessário para a avaliação.

\section{Considerações finais do capítulo}

Este capítulo apresentou a metodologia proposta para estimação de fasores harmônicos em tempo real. O método combina rastreamento da frequência fundamental por Zero-Crossing, reamostragem dinâmica por interpolação B-Spline e extração harmônica por banco de filtros Flat-Top em estrutura polifásica. A principal ideia é adaptar a base temporal do sinal antes da análise espectral, reduzindo o impacto de desvios da frequência fundamental sobre os fasores harmônicos estimados.

As etapas descritas também estabelecem a ponte entre o modelo em software e a implementação embarcada em desenvolvimento. A descrição foi mantida no nível metodológico, sem assumir validação final em FPGA ou resultados conclusivos. A fundamentação teórica necessária para compreender os blocos utilizados é apresentada no capítulo anterior, enquanto a análise de resultados será incorporada posteriormente, após a consolidação dos ensaios.
